% curly.tex version 1.01, last modified 12 July 1993 by A. Lewenberg.
%
% This file defines the Formal Script family of fonts based on the 
% Formal Script Math Symbols 10 point fonts by Ralph A. Smith. 
% It uses a \newfam, so use with caution. It is compatible with PLAIN
% and AMSTEX 2.1: in PLAIN it functions in the same way as \cal; in
% AMSTEX it functions the same way as \Cal. Note that, just like \cal
% and \Cal, the command \curly only works in math mode.
%
% To use, put an `\input curly` somewhere in your file.
% The command \curly changes to the curly script family. For
% example, $a={\curly a}$.
% Note: You must have the fonts rsfs10, rsfs7 and rsfs5.
%                Adam H. Lewenberg adam@symcom.math.uiuc.edu
%
%%%%%%%%%%%%%%%%%%%%%%%%%%%%%%%%%%%%%%%%%%%%%%%%%%%%%%%%%%%%%%%%%%%%%%
%
% Because the following code uses a \newfam,
% make sure to load the file only once. 
%
\ifx\curlyloaded\relax\let\next\endinput
  \else\let\curlyloaded\relax\let\next\relax
\fi
\next
%
% Save the catcode of @
\xdef\recoveratcodezqrz{\catcode`\noexpand\@=\the\catcode`\@}
\catcode`\@=11
%
\edef\innernewfam{\expandafter\noexpand\csname newfam\endcsname}
%
% Define the font family. Do different things if AMSTEX loaded. 
\ifx\amstexloaded@\relax
% AMSTEX loaded
  \loadmathfont{rsfs}
  \def\curly{\mathfont@\curly}
  \def\curlyfam{\rsfsfam}
\else
% AMSTEX not loaded; assume PLAIN and act like \cal
  \font\tenrsfs=rsfs10
  \font\sevenrsfs=rsfs7
  \font\fiversfs=rsfs5
  \innernewfam\rsfsfam \def\curly{\fam\rsfsfam\tenrsfs}
  \textfont\rsfsfam=\tenrsfs  \scriptfont\rsfsfam=\sevenrsfs 
  \scriptscriptfont\rsfsfam=\fiversfs
  \def\curly{\fam\rsfsfam}
\fi
\let\innernewfam\undefined
\recoveratcodezqrz
\endinput%
\endinput%
